% ===================================================================== METHODS
% Target: about 1.5 columns. Describe what was done so another student could repeat it.
% Every number below is a \val{} macro; keep using them instead of typing digits.
\section{Methods}
\label{sec:methods}

% --- A. Data. two_curves from the benchmark suite of \cite{bowles2024benchmarking};
% breast_cancer (UCI) reduced to \val{n-qubits} features with PCA fitted on the training split.
% Features scaled to [0, pi] using training-split ranges. Fixed stratified splits:
% train \val{n-train}, validation \val{n-val}, labelled hardware calibration pool \val{n-cal},
% hardware test set \val{n-test} (the same samples in every session). Seed \val{seed}.
\subsection{Datasets}
\todo{Datasets and splits.}

% --- B. Models. Three circuits on \val{n-qubits} qubits, each followed by a logistic head on
% Pauli expectation features (state the feature count for each):
%  M1 data re-uploading \cite{perezsalinas2020reuploading}: \val{m1-layers} layers of RY(x),
%     trainable RZ RY per qubit and a CZ chain, plus a final RY layer; \val{m1-params}
%     parameters trained jointly with the head (Adam, exact parameter shift, \val{m1-steps} steps,
%     \val{m1-restarts} restarts); features <Z_q> (\val{m1-features}).
%  M2 projected features \cite{huang2021power}: \val{m2-layers} fixed layer of RY(x) RZ(x) and a
%     CZ chain; features <X_q>,<Y_q>,<Z_q> (\val{m2-features}), three measurement circuits.
%  M3 quantum extreme learning machine: data layer, \val{m3-layers} fixed random layers;
%     features <Z_q> and <Z_q Z_r> (\val{m3-features}).
%  Head: p(y=1|x) = sigma(w^T phi(x) + b) on standardised features, C chosen by 5-fold CV.
%  Useful equation:
%   p(y=1\mid x) = \sigma\!\left(w^\top \tfrac{\phi(x)-\mu}{s} + b\right),\quad
%   \phi_j(x) = \langle P_j \rangle_{\rho(x)}
%  Say that M2/M3 depth was chosen on validation accuracy before any hardware run.
\subsection{Models}
\todo{Models M1--M3 and the classical head.}

% --- C. Hardware protocol. Heron r2 devices \val{devices} (fallback \val{fallback}); Open Plan,
% one Sampler job per session, \val{shots} shots per circuit, \val{readout-shots} shots for the
% two readout-calibration circuits; twirling and dynamical decoupling off (raw noise).
% Packing: \val{patches} disjoint linear \val{n-qubits}-qubit patches with \val{patch-buffer}
% idle qubit between them, chosen once from calibration data and frozen; hardware sample i runs
% on patch i mod \val{patches} in every session. Cite \cite{parallelqml2024,ohkura2022simultaneous}.
% Pilot: packed vs isolated (Table~\ref{tab:pilot}); measured cost \val{per-execution-us}~$\mu$s
% per execution. Sessions: \val{sessions-per-device} per device over 10-14 days (Table~\ref{tab:protocol}).
% Budget: hard cap \val{budget-cap}~s across accounts; campaign cap \val{campaign-cap}~s;
% \val{qpu-used}~s used so far. Calibration snapshots saved at submission and at job start.
\subsection{Hardware protocol}
\todo{Devices, packing, pilot, sessions and budget.}

% --- D. Emulator and repairs. EMU: AerSimulator.from_backend with the calibration available
% at submission (\val{emulator-shots} shots). Repairs, all computed offline from the same counts:
%  R1 batch normalisation of features \cite{wang2022quantumnat};
%  R2_k head refit on k = 8, 16, 32 labelled hardware samples, penalised toward the simulator
%     head with strength \val{refit-lambda} k; R2_16^s1 reuses session 1 labels;
%  R3 tensored readout mitigation with measured (R3m) or reported (R3r) error rates \cite{nation2021m3};
%  R4 head trained on emulated noisy training features (\val{noise-aware-shots} shots);
%  extra QPU seconds counted for each repair.
\subsection{Emulator and repairs}
\todo{Emulator and repairs R1--R4.}

% --- E. Uncertainty. Split conformal with score 1 - p(true label) calibrated on the simulator
% (CP_sim), on session 1 (CP_s1) or on the same session (CP_same) \cite{park2023qcp}; adaptive
% conformal inference with step \val{aci-gamma} over the session stream \cite{gibbs2021aci,
% spencer2026aqcp}; target coverage \val{coverage-target}. Calibration error ECE \cite{guo2017calibration}.
\subsection{Uncertainty quantification}
\todo{Conformal prediction and calibration metrics.}

% --- F. Statistics and pre-registration. Hypotheses H1-H5 and decision rules fixed in
% preregistration.md and tagged prereg-v1 before session 1; \val{bootstrap} bootstrap
% resamples; OLS with HC3 errors for H4; list deviations (notes/deviations.md).
\subsection{Statistics and pre-registration}
\todo{Hypotheses, tests and deviations.}

\begin{table}[t]
\centering
\caption{Noiseless simulator test accuracy of the quantum models and tuned classical baselines (150 test samples; 95\% bootstrap interval).}
\label{tab:sim}
\footnotesize
\setlength{\tabcolsep}{3.5pt}
\begin{tabular}{lcc}
\toprule
Model & Two curves & Breast cancer \\
\midrule
M1 re-uploading & 0.927 [0.880, 0.967] & 0.967 [0.933, 0.993] \\
M2 projected & 0.873 [0.820, 0.927] & 0.933 [0.893, 0.973] \\
M3 QELM & 0.847 [0.787, 0.900] & 0.913 [0.867, 0.960] \\
\midrule\multicolumn{3}{l}{\emph{Tuned classical baselines, same inputs}} \\
Logistic regression & 0.733 [0.660, 0.800] & 0.940 [0.900, 0.973] \\
RBF-SVM & 0.987 [0.967, 1.000] & 0.940 [0.900, 0.973] \\
MLP & 0.973 [0.947, 0.993] & 0.960 [0.927, 0.987] \\
\bottomrule
\end{tabular}
\end{table}

\begin{table}[t]
\centering
\caption{Hardware runs on \texttt{ibm\_}* devices: job start (2026, UTC), queue wait, billed QPU time and age of the reported calibration at the job start.}
\label{tab:protocol}
\footnotesize
\setlength{\tabcolsep}{3.5pt}
\begin{tabular}{llrrrr}
\toprule
Device & Run & Start & Queue (h) & QPU (s) & Cal.\ age (h) \\
\midrule
fez & pilot & 10-26 15:33 & 0.0 & 11 & 0.4 \\
fez & 1 & 10-27 03:37 & 0.0 & 20 & 0.3 \\
fez & 2 & 10-30 18:24 & 0.3 & 20 & 0.3 \\
fez & 3 & 11-04 05:07 & 0.3 & 20 & 0.2 \\
fez & 4 & 11-10 17:09 & 0.0 & 20 & 0.6 \\
kingston & pilot & 10-26 21:38 & 4.9 & 11 & 0.2 \\
kingston & 1 & 10-27 05:31 & 1.9 & 21 & 0.3 \\
kingston & 2 & 10-30 18:23 & 0.2 & 21 & 0.4 \\
kingston & 3 & 11-04 04:47 & 0.0 & 21 & 0.2 \\
kingston & 4 & 11-10 17:09 & 0.0 & 21 & 0.8 \\
kingston & 5 & 11-10 20:04 & 0.0 & 21 & 0.1 \\
\bottomrule
\end{tabular}
\end{table}

\begin{table}[t]
\centering
\caption{Pilot: the same samples run packed with all patches active or alone on their patch. Excess total-variation distance over a shot-noise null, and accuracy difference (packed minus isolated), with 95\% bootstrap intervals.}
\label{tab:pilot}
\footnotesize
\setlength{\tabcolsep}{3.5pt}
\begin{tabular}{llcccc}
\toprule
Device & Model & Excess TVD & Packed & Isolated & $\Delta$acc \\
\midrule
fez & M1 & [-0.002, +0.019] & 0.82 & 0.79 & [-0.07, +0.14] \\
fez & M3 & [-0.009, +0.006] & 0.64 & 0.68 & [-0.21, +0.14] \\
kingston & M1 & [-0.013, +0.004] & 0.79 & 0.82 & [-0.18, +0.07] \\
kingston & M3 & [-0.012, +0.004] & 0.71 & 0.71 & [-0.18, +0.18] \\
\bottomrule
\end{tabular}
\end{table}

